\documentclass[10pt,a4paper]{amsart}
\usepackage[margin=19mm]{geometry}
\usepackage{amsmath,amssymb,mathtools}
\usepackage[scr=rsfs,cal=euler]{mathalfa}
\usepackage{xfrac}
\usepackage[unicode,hidelinks,bookmarks=false]{hyperref}
\newcommand{\ie}{i.e.\ }
\newcommand{\cf}{cf.\ }
\newcommand{\resp}{respectively\ }
\newcommand{\Subsection}{\subsection}
\providecommand{\nobreakdash}{\nobreak\hbox{-}}
\setlength{\emergencystretch}{2em}
\multlinegap0pt
\usepackage{amssymb}
\usepackage{mathtools}
\usepackage[shortlabels]{enumitem}
\newtheorem{proposition}{Proposition}
\newcommand{\Z}{\mathbb{Z}}
\title{On the distribution of shapes of octic Kummer extensions}
\author{Anuj Jakhar, Anwesh Ray}
\date{Source: arXiv:2602.12621v1 (CC BY 4.0)}
\begin{document}
\maketitle

\noindent\emph{Source and licence.} On the distribution of shapes of octic Kummer extensions. Version arXiv:2602.12621v1 (CC BY 4.0), distributed by arXiv under the Creative Commons Attribution 4.0 International licence, http://creativecommons.org/licenses/by/4.0/, as recorded in the arXiv licence metadata for this exact version and shown on the abstract page https://arxiv.org/abs/2602.12621v1. Full licence notice: \texttt{licenses/arxiv-2602.12621-CC-BY-4.0.txt}.\par\smallskip

\noindent\textit{Editorial scope.} Counted unit: the article's proposition count (source0.tex lines 1488--1500). The article's complete original proof of it is delivered with it (source0.tex lines 1502--1708, one \texttt{proof} environment of 207 lines). No mathematics is written, repaired or re-derived: the statement and the proof are the article's own lines, in the article's own notation, and no retained line is substituted or re-flowed.  Retained uncounted as the source-local context the proof needs: lines 1473--1486.  Nothing was omitted from any retained span, so the package carries no omission record.  No retained line contains a citation, so the wrapper carries no bibliography and no bibliography entry.  No retained line contains a figure environment, an image-inclusion command or drawing code, so no image is delivered and the package declares no asset dependency.  Editorial formatting only, in addition: an AMS \texttt{amsart} preamble that restates the article's own declarations and macros used by the retained lines, namely the environments \texttt{proposition} on separate counters, together with the standard packages those lines need.  The retained environments print in this document as Proposition 1 in order of appearance, whereas the article prints the same items with the numbering of its own full document.  The complete native source of the article as distributed in the arXiv e-print for this exact version is bundled unchanged as \texttt{sources/194566/source0.tex}.
\begin{equation}\label{eq:Nprime-sum-reduction}
\begin{aligned}
N'(\mathcal R;X)
&=
\sum_{v}N_v\sum_{w}
\pi N_w\left(
\min\!\left(R_1w,\frac{X}{v^2w^3}\right)^2
-
(R_1'w)^2
\right)
+
O_{\mathcal R}(X^{1/2+\varepsilon}).
\end{aligned}
\end{equation}

\begin{proposition}\label{Propn N' count}
As $X\to\infty$, one has
\[
N'(\mathcal R;X)
=
\pi^2\bigl(R_1-R_1'\bigr)
X
\sum_{v\in[R_2',R_2]}
\frac{N_v}{v^2}
+
O_{\mathcal R}(X^{3/4}).
\]
\end{proposition}

\begin{proof}
We begin from \eqref{eq:Nprime-sum-reduction}:
\[
N'(\mathcal R;X)
=
\pi
\sum_{v} N_v
\sum_{w}
N_w
\Big(
\min(R_1w,X/(v^2w^3))^2
-
(R_1'w)^2
\Big)
+
O_{\mathcal R}(X^{1/2+\varepsilon}),
\]
where $v\in[R_2',R_2]$ and the inner sum runs over $w$ for which the
summand is nonzero.

Fix $v$ in this interval.  Define
\[
F(t)
:=
\min(R_1t,X/(v^2t^3))^2
-
(R_1't)^2,
\]
and let
\[
T_v := \left(\frac{X}{R_1'v^2}\right)^{1/4}.
\]
Then the inner sum may be written as
\[
\sum_{w} N_w F(w)
=
\sum_{\substack{h\in\Z[i]\\ |h|\le T_v}} F(|h|).
\]
Setting $A(T):=\#\{h\in\Z[i]: |h|\le T\}$, one has that
\[
A(T)=\pi T^2+O(T).
\]
\noindent Write $0=r_0<r_1<\cdots<r_N\le T_v$ for the distinct values of $|h|$ in
increasing order.  Then
\[
\sum_{\substack{h\in\Z[i]\\ |h|\le T_v}} F(|h|)
=
\sum_{j=1}^N F(r_j)\bigl(A(r_j)-A(r_{j-1})\bigr).
\]
\noindent Applying Abel summation gives
\[
=
F(r_N)A(r_N)
-
\sum_{j=1}^{N-1} A(r_j)\bigl(F(r_{j+1})-F(r_j)\bigr).
\]
Since $F(t)\ll 1$ and $A(r_N)\ll T_v^2\asymp X^{1/2}$, we obtain
\[
F(r_N)A(r_N)=O(X^{1/2}).
\]
By the mean value theorem,
\[
F(r_{j+1})-F(r_j)
=
F'(\xi_j)(r_{j+1}-r_j)
\]
for some $\xi_j\in(r_j,r_{j+1})$.  Therefore
\[
\sum_{\substack{h\in\Z[i]\\ |h|\le T_v}} F(|h|)
=
-\sum_{j} A(r_j)F'(\xi_j)(r_{j+1}-r_j)
+
O(X^{1/2}).
\]
\noindent The sum on the right is a Riemann sum for the integral
\[
-\int_0^{T_v} A(r)F'(r)\,dr.
\]
Since $A(r)=\pi r^2+O(r)$, the difference between the discrete sum and the
integral is bounded by
\[
\ll \int_0^{T_v} r\,|F'(r)|\,dr.
\]
Hence
\[
\sum_{\substack{h\in\Z[i]\\ |h|\le T_v}} F(|h|)
=
-\int_0^{T_v} A(r)F'(r)\,dr
+
O\!\left(
\int_0^{T_v} r\,|F'(r)|\,dr
+
X^{1/2}
\right).
\]
\noindent Substituting $A(r)=\pi r^2+O(r)$ gives
\[
=
-\pi\int_0^{T_v} r^2F'(r)\,dr
+
O\!\left(
\int_0^{T_v} r\,|F'(r)|\,dr
+
X^{1/2}
\right).
\]
Note that $F(T_v)=F(0)=0$. Integration by parts yields
\[
-\pi\int_0^{T_v} r^2F'(r)\,dr
=
2\pi\int_0^{T_v} rF(r)\,dr.
\]

We may now evaluate these integrals. Let 
\[
r_0 := \left( \frac{X}{R_1 v^2} \right)^{1/4}
\quad\text{and}\quad
r_1 := \left( \frac{X}{R_1' v^2} \right)^{1/4}.
\]
By definition, the function $F$ is given by:
\[
F(r) = 
\begin{cases} 
(R_1^2 - R_1'^2) r^2, & 0 \le r \le r_0,\\[2mm]
\dfrac{X^2}{v^4 r^6} - R_1'^2 r^2, & r_0 < r \le r_1,\\[1mm]
0, & r > r_1.
\end{cases}
\]
First, consider the integral
\[
\int_0^{T_v} r F(r)\, dr 
= \int_0^{r_0} r (R_1^2 - R_1'^2) r^2\, dr 
+ \int_{r_0}^{r_1} r \left( \frac{X^2}{v^4 r^6} - R_1'^2 r^2 \right) dr.
\]
We have
\[
\int_0^{r_0} (R_1^2 - R_1'^2) r^3\, dr = \frac{R_1^2 - R_1'^2}{4} r_0^4 = \frac{R_1^2 - R_1'^2}{4} \cdot \frac{X}{R_1 v^2} = \frac{X (R_1 - R_1')(R_1 + R_1')}{4 R_1 v^2},
\]
and
\[
\int_{r_0}^{r_1} r \left( \frac{X^2}{v^4 r^6} - R_1'^2 r^2 \right) dr
= \frac{X^2}{v^4} \int_{r_0}^{r_1} r^{-5} dr - R_1'^2 \int_{r_0}^{r_1} r^3 dr.
\]

We compute each term separately:

\[
\frac{X^2}{v^4} \int_{r_0}^{r_1} r^{-5} dr = \frac{X^2}{v^4} \cdot \frac{r_0^{-4} - r_1^{-4}}{4} 
= \frac{X^2}{4 v^4} \left( \frac{R_1 v^2}{X} - \frac{R_1' v^2}{X} \right) 
= \frac{X (R_1 - R_1')}{4 v^2},
\]

and

\[
R_1'^2 \int_{r_0}^{r_1} r^3 dr = \frac{R_1'^2}{4} (r_1^4 - r_0^4) 
= \frac{R_1'^2}{4} \cdot \frac{X}{v^2} \left( \frac{1}{R_1'} - \frac{1}{R_1} \right) 
= \frac{X (R_1 - R_1') R_1'}{4 v^2 R_1}.
\]

Hence the second integral simplifies to
\[
\int_{r_0}^{r_1} r F(r)\, dr = \frac{X (R_1 - R_1')}{4 v^2} - \frac{X (R_1 - R_1') R_1'}{4 v^2 R_1} 
= \frac{X (R_1 - R_1')^2}{4 R_1 v^2}.
\]
\noindent Combining the two contributions, we obtain
\[
\int_0^{T_v} r F(r)\, dr = \frac{X (R_1 - R_1')(R_1 + R_1')}{4 R_1 v^2} + \frac{X (R_1 - R_1')^2}{4 R_1 v^2} = \frac{X (R_1 - R_1')}{2 v^2}.
\]
\noindent Multiplying by \(2\pi\) to account for angular directions gives
\[
2\pi \int_0^{T_v} r F(r)\, dr = \frac{\pi X (R_1 - R_1')}{v^2}.
\]

For the derivative, we have
\[
F'(r) = 
\begin{cases} 
2 (R_1^2 - R_1'^2) r, & 0 \le r \le r_0,\\
-6 \frac{X^2}{v^4 r^7} - 2 R_1'^2 r, & r_0 < r \le r_1.
\end{cases}
\]
Hence
\[
\int_0^{T_v} r |F'(r)| dr \ll \int_0^{r_0} r^2 dr + \int_{r_0}^{r_1} \left( \frac{X^2}{v^4 r^6} + r^2 \right) dr \ll \frac{X^{3/4}}{v^{3/2}}.
\]
\noindent Since \(v\) lies in a fixed interval \([R_2', R_2]\), this gives an error term
\[
O_{\mathcal R}(X^{3/4}).
\]
We conclude that
\[
\sum_{|h|\leq T_v} F(|h|) = \frac{\pi X (R_1 - R_1')}{v^2} + O_{\mathcal R}(X^{3/4}).
\]
Substitute this into the outer sum over $v$ with $\varepsilon\in (0, 1/4)$ to get:
\[
N'(\mathcal R;X)
=
\pi^2\bigl(R_1-R_1'\bigr)
X
\sum_{v\in[R_2',R_2]}
\frac{N_v}{v^2}
+
O_{\mathcal R}(X^{3/4}),
\]
which completes the proof.
\end{proof}

\end{document}
